%% maestro2tex -- generated figure; do not edit by hand.
\providecommand{\MaestroRoot}{include/analog/}
\begin{figure}[htbp]
  \centering
  {\pgfplotsset{compat=1.18}%
  \begin{tikzpicture}
    \begin{axis}[
      width=0.82\linewidth, height=6.2cm,
      xlabel={Frequency~[\si{\hertz}]}, ylabel={Open-loop gain $|A|$~[\si{\decibel}]},
      grid=both,
      major grid style={line width=0.2pt, draw=black!14},
      minor grid style={line width=0.2pt, draw=black!7},
      minor tick num=1,
      axis line style={line width=0.4pt, draw=black!45},
      tick style={line width=0.4pt, draw=black!45},
      tick align=outside,
      label style={font=\small}, tick label style={font=\small},
      enlarge x limits=0.02, enlarge y limits=0.10,
      every axis plot/.append style={line join=round, no marks},
      legend style={font=\footnotesize, draw=black!25, fill=white,
                    fill opacity=0.9, text opacity=1, row sep=1pt},
      legend cell align=left, legend pos=south west,
      xmode=log,
      scaled y ticks=false,
      y tick label style={/pgf/number format/fixed,
                          /pgf/number format/precision=0}
    ]
      \addplot[mtxA, line width=1.1pt, solid] table {\MaestroRoot data/tia_loopgain_00.dat};
      \addlegendentry{tt}
      \addplot[mtxB, line width=1.1pt, dashed] table {\MaestroRoot data/tia_loopgain_01.dat};
      \addlegendentry{ff}
      \addplot[mtxC, line width=1.1pt, dotted] table {\MaestroRoot data/tia_loopgain_02.dat};
      \addlegendentry{fs}
      \addplot[mtxD, line width=1.1pt, dashdotted] table {\MaestroRoot data/tia_loopgain_03.dat};
      \addlegendentry{sf}
      \addplot[mtxE, line width=1.1pt, densely dashed] table {\MaestroRoot data/tia_loopgain_04.dat};
      \addlegendentry{ss}
    \end{axis}
  \end{tikzpicture}}
  \caption[{Open-loop gain of the transimpedance amplifier over the five process corners, by \texttt{stb} analysis\,\dots}]{Open-loop gain of the transimpedance amplifier over the five process corners, by \texttt{stb} analysis through the feedback probe with the amplifier in unity feedback and a \SI{5}{\pico\farad} load. Miller compensation places the dominant pole near \SI{200}{\hertz} and gives a single-pole roll-off to a \SIrange{6.2}{7.6}{\mega\hertz} crossover; the corners overlay to within \SI{1.7}{\decibel} of DC gain. This is the amplifier's own response. The loop that closes around the electrode is a different loop, and it is not this well behaved --- see Figure~\ref{fig:tia-zt}.}
  \label{fig:tia-loopgain}
\end{figure}
